% !TEX program = lualatex
\documentclass[aspectratio=43,professionalfonts,handout]{beamer}
\usepackage{beamer_template}

\newcommand{\LectureNum}{13}
\newcommand{\LectureTitle}{フーリエ変換の性質と公式}
\newcommand{\TextPages}{p.\,96-99}
\newcommand{\CourseName}{応用数学}
\renewcommand{\TermName}{前期}

\begin{document}

\begin{frame}{\CourseName\quad 第\LectureNum 回「\LectureTitle」}
  {\small \TermName\quad \TeacherName\quad ／\quad 教科書 \TextPages}

  \textbf{目標}：フーリエ変換の基本的性質を理解し、たたみ込みのフーリエ変換を計算できる

\end{frame}

\begin{frame}{フーリエ変換の性質一覧}
  \begin{block}{}
  \begin{itemize}
    \item (I) 線形性
    \item (II) 推移：$e^{-iau}F(u)$
    \item (III) 相似性：$(1/|a|)F(u/a)$
    \item (IV) 微分：$(iu)^n F(u)$
    \item (V)(VI) 高次微分
  \end{itemize}
  \end{block}
\end{frame}

\begin{frame}{たたみ込みのフーリエ変換}
  \begin{block}{}
  \begin{itemize}
    \item $(f*g)(x) = \int_{-\infty}^{\infty} f(x-\xi)g(\xi)d\xi$
    \item $\mathcal{F}[f*g] = F(u)G(u)$
    \item ラプラス変換との対応
  \end{itemize}
  \end{block}
\end{frame}

\begin{frame}{ガウス関数}
  \begin{block}{}
  \begin{itemize}
    \item $\mathcal{F}[e^{-ax^{2}}] = \sqrt{\pi/a}e^{-u^{2}/4a}$
    \item $a=1/2$ のとき $\mathcal{F}[e^{-x^{2}/2}] = \sqrt{2\pi}e^{-u^{2}/2}$
    \item 確率・統計との関係
  \end{itemize}
  \end{block}
\end{frame}

\begin{frame}{演習}
  \begin{exampleblock}{自分で解いてみよう}
  \begin{itemize}
    \item 問5〜9より選題
  \end{itemize}
  \end{exampleblock}
\end{frame}

\begin{frame}{まとめ}
  \begin{itemize}
    \item フーリエ変換の6つの性質
    \item たたみ込みとの関係
    \item 次回：スペクトル
  \end{itemize}
\end{frame}

\end{document}
